% meshgraphnet_slides.tex -- additional MeshGraphNet slides
% Source context: MIKU-7437/AI_Termoelastic_Waves_Geology, MeshGraphNet files.

\begin{frame}{MeshGraphNet: Why a Graph Model}
  \begin{columns}[T]
    \begin{column}{0.48\textwidth}
      \begin{block}{Reason for using MeshGraphNet}
        \footnotesize
        COMSOL finite-element data are naturally mesh-based. Nodes can have
        different neighbourhoods and nonuniform distances, unlike pixels on a
        regular image grid.
      \end{block}
      \begin{itemize}
        \footnotesize
        \item Preserves unstructured mesh connectivity.
        \item Avoids forced interpolation to a Cartesian grid.
        \item Matches local thermal and mechanical interactions.
      \end{itemize}
    \end{column}
    \begin{column}{0.48\textwidth}
      \begin{block}{Graph representation}
        \[
          \mathcal{G}=(\mathcal{V},\mathcal{E})
        \]
      \end{block}
      \scriptsize
      $\mathcal{V}$ -- mesh nodes; $\mathcal{E}$ -- local graph edges.
      Each node represents a spatial location in the simulated geological
      medium; each edge represents local physical neighbourhood connectivity.
    \end{column}
  \end{columns}
\end{frame}

\begin{frame}{Mesh and Edge Features}
  \begin{columns}[T]
    \begin{column}{0.50\textwidth}
      \begin{block}{Geometric edge encoding}
        \[
          r_{ij}=p_j-p_i,\qquad d_{ij}=\|r_{ij}\|_2
        \]
        \[
          a_{ij}=\begin{bmatrix}r_{ij}\\ d_{ij}\end{bmatrix}
        \]
      \end{block}
    \end{column}
    \begin{column}{0.46\textwidth}
      \footnotesize
      The edge attributes describe both direction and distance between
      neighbouring mesh nodes.
      \vskip 4pt
      This matters because heat transfer, stress redistribution, and elastic
      wave response depend on local geometry, not only on node values.
      \vskip 4pt
      \scriptsize
      $p_i,p_j$ -- node coordinates [m]; $r_{ij}$ -- relative vector [m];
      $d_{ij}$ -- Euclidean edge length [m].
    \end{column}
  \end{columns}
\end{frame}

\begin{frame}{Node Features and Prediction Target}
  \begin{columns}[T]
    \begin{column}{0.48\textwidth}
      \begin{block}{Node feature vector}
        \[
          x_i^t=\begin{bmatrix}p_i\\ m_i\\ s\\ q_i^t\end{bmatrix}
        \]
      \end{block}
      \footnotesize
      Coordinates, material parameters, scenario conditions, and current
      dynamic fields are concatenated before model input.
    \end{column}
    \begin{column}{0.48\textwidth}
      \begin{block}{State transition}
        \[
          \hat{Y}^{t+\Delta t}=f_{\theta}(X^t,\mathcal{G})
        \]
      \end{block}
      \scriptsize
      $X^t\in\mathbb{R}^{N\times F_v}$ is the node feature matrix;
      $\hat{Y}^{t+\Delta t}\in\mathbb{R}^{N\times C}$ is the predicted
      future field matrix for all graph nodes.
    \end{column}
  \end{columns}
\end{frame}

\begin{frame}{Delta Prediction Mode}
  \begin{columns}[T]
    \begin{column}{0.48\textwidth}
      \begin{block}{Absolute mode}
        \[
          \hat{q}_i^{t+\Delta t}=f_{\theta}(x_i^t,\mathcal{G})
        \]
      \end{block}
      \begin{block}{Delta mode}
        \[
          \Delta\hat{q}_i^t=f_{\theta}(x_i^t,\mathcal{G})
        \]
        \[
          \hat{q}_i^{t+\Delta t}=q_i^t+\Delta\hat{q}_i^t
        \]
      \end{block}
    \end{column}
    \begin{column}{0.48\textwidth}
      \footnotesize
      Delta prediction is useful because neighbouring time-step changes are
      often smoother and smaller than full thermoelastic fields.
      \vskip 5pt
      \begin{alertblock}{Rollout caveat}
        \footnotesize
        In autoregressive prediction, each predicted state becomes the next
        input. Small one-step errors can accumulate over long horizons.
      \end{alertblock}
    \end{column}
  \end{columns}
\end{frame}

\begin{frame}{Encoder--Processor--Decoder Structure}
  \begin{block}{Model composition}
    \[
      f_{\theta}=D_{\theta_D}\circ P_{\theta_P}\circ E_{\theta_E}
    \]
  \end{block}
  \vskip 4pt
  \begin{columns}[T]
    \begin{column}{0.31\textwidth}
      \begin{block}{Encoder}
        \footnotesize
        Maps raw node and edge features into latent vectors.
      \end{block}
    \end{column}
    \begin{column}{0.31\textwidth}
      \begin{block}{Processor}
        \footnotesize
        Applies repeated message passing over mesh connectivity.
      \end{block}
    \end{column}
    \begin{column}{0.31\textwidth}
      \begin{block}{Decoder}
        \footnotesize
        Converts final latent node states into predicted fields.
      \end{block}
    \end{column}
  \end{columns}
  \vskip 4pt
  \scriptsize
  The repository implementation uses multilayer perceptrons, smooth SiLU
  activations, residual connections, and layer normalization inside the graph
  processing path.
\end{frame}

\begin{frame}{Message Passing, Training, and Limits}
  \begin{columns}[T]
    \begin{column}{0.52\textwidth}
      \begin{block}{Message passing step}
        \[
          \tilde{e}_{ij}^{k+1}=\phi_e^k([h_i^k,h_j^k,e_{ij}^k])
        \]
        \[
          m_j^{k+1}=\sum_{i:(i,j)\in\mathcal{E}}e_{ij}^{k+1}
        \]
        \[
          h_j^{k+1}=\phi_v^k([h_j^k,m_j^{k+1}])
        \]
      \end{block}
    \end{column}
    \begin{column}{0.44\textwidth}
      \begin{block}{Training objective}
        \[
          \mathcal{L}=\frac{1}{\sum_c w_c}\sum_{c=1}^{C}w_c
          \frac{1}{N}\sum_{i=1}^{N}(\hat{y}_{i,c}-y_{i,c})^2
        \]
      \end{block}
      \scriptsize
      MeshGraphNet is a supervised COMSOL surrogate. It preserves mesh topology
      and supports material conditioning, but physical equations are learned
      implicitly and must be checked through validation plots and metrics.
    \end{column}
  \end{columns}
\end{frame}
